\section{Algorithm for estimating ITE}\label{sec:model}

We propose a general framework called CFR (for Counterfactual Regression) for ITE estimation based on the theoretical results above. Our algorithm is an end-to-end, regularized minimization procedure which simultaneously fits both a balanced representation of the data and a hypothesis for the outcome.
CFR draws on the same intuition as the approach proposed by  \citet{johansson2016counterfactual}, but \emph{overcomes} the following limitations of their method: a) Their theory requires a two-step optimization procedure and is specific to \emph{linear} hypotheses of the learned representation (and does not support e.g. deep neural networks), b) The treatment indicator might get lost if the learned representation is high-dimensional (see discussion below).

%We note that our framework is also more flexible in practice, as our theory supports multiple measures of balance that can can be minimized efficiently; this is only rarely true for variants of the discrepancy distance used by \cite{johansson2016counterfactual}.

We assume there exists a distribution $p(x,t,Y_0,Y_1)$ over $\cX \times \{0,1\} \times \cY \times \cY$, such that strong ignorability holds. We further assume we have a sample from that distribution $(x_1,t_1,y_1), \dots (x_n,t_n,y_n)$, where $y_i \sim p(Y_1|x_i)$ if $t_i=1$, $y_i \sim p(Y_0|x_i)$ if $t_i=0$. This standard assumption means that the treatment assignment determines which potential outcome we see. Our goal is to find a representation $\Phi : \cX \rightarrow \cR$ and hypothesis $h:  \cX \times \{0,1\} \rightarrow \cY$ that will minimize $\epehe(f)$ for $f(x,t):=h(\Phi(x),t)$.

In this work, we let $\Phi(x)$ and $h(\Phi,t)$ be parameterized by deep neural networks trained jointly in an end-to-end fashion, see Figure~\ref{fig:neuralnet}. This model allows for learning complex non-linear representations and hypotheses with large flexibility. \citet{johansson2016counterfactual} parameterized $h(\Phi,t)$ with a single network using the concatenation of $\Phi$ and $t$ as input. When the dimension of $\Phi$ is high, this risks losing the influence of $t$ on $h$ during training. To combat this, our first contribution is to parameterize $h_1(\Phi)$ and $h_0(\Phi)$ as two separate ``heads'' of the joint network, the former used to estimate the outcome under treatment, and the latter under control. This means that statistical power is shared in the representation layers of the network, while the effect of treatment is retained in the separate heads. Note that each sample is used to update only the head corresponding to the observed treatment; for example, an observation $(x_i, t_i=1, y_i)$ is only used to update $h_1$.

Our second contribution is to excplicitly account and adjust for the bias induced by treatment group imbalance. To this end, we seek a representation $\Phi$ and hypothesis $h$ that minimizes a trade-off between predictive accuracy and imbalance in the representation space, using the following objective:
\begin{equation}
\begin{aligned}
\min_{\substack{h, \Phi\\ \|\Phi\|=1}} & \;\; \frac 1 n \sum_{i=1}^n w_i \cdot L\left(h(\Phi(x_i),t_i)\, ,y_i\right) + \lambda \cdot \mathfrak{R}(h)  \\
& \;\; + \alpha  \cdot \text{IPM}_\cF\left(\{\Phi(x_i)\}_{i:t_i=0},\{\Phi(x_i)\}_{i:t_i=1}\right), \\
\mbox{with} & \;\; w_i = \frac{t_i}{2u} + \frac{1-t_i}{2(1-u)},\;\;\mbox{where}\;\;u = \frac{1}{n}\sum_{i=1}^n t_i, \\
\mbox{and} & \;\;\mbox{$\mathfrak{R}$ is a model complexity term.}
\label{eq:overallalg}
\end{aligned}
\end{equation}
Note that $u=p(t=1)$ in the definition of $w_i$ is simply the proportion of treated units in the population. The weights $w_i$ compensate for the difference in treatment group size in our sample, see Theorem~\ref{thm:gen}. $\text{IPM}_{\cF}(\cdot,\cdot)$ is the (empirical) integral probability metric defined by the function family $\cF$. For most IPMs, we cannot compute the factor $B_\phi$ in Equation~\ref{eq:thm_gen}, but treat it as part of the hyperparameter $\alpha$. This makes our objective sensitive to the scaling of $\Phi$, even for a constant $\alpha$. We therefore normalize $\Phi$ through either projection or batch-normalization with fixed scale. We refer to the model minimizing \eqref{eq:overallalg} with $\alpha>0$ as Counterfactual Regression (CFR) and the variant without balance regularization ($\alpha=0$) as Treatment-Agnostic Representation Network (\tarnet{}).

We train our models by minimizing \eqref{eq:overallalg} using stochastic gradient descent, where we backpropagate the error through both the hypothesis and representation networks, as described in Algorithm~\ref{alg:model}. Both the prediction loss and the penalty term $\text{IPM}_{\cF}(\cdot,\cdot)$ are computed for one mini-batch at a time. Details of how to obtain the gradient $g_1$ with respect to the empirical IPMs are in the supplement.

\begin{algorithm}[tbp]
\caption{CFR: Counterfactual regression with integral probability metrics}
\label{alg:model}
\begin{algorithmic}[1]
  \STATE \textbf{Input:} Factual sample $(x_1,t_1,y_1), \ldots , (x_n,t_n,y_n)$, scaling parameter $\alpha>0$, loss function $L\left(\cdot,\cdot\right)$, representation network $\Phi_{\bf{W}}$ with initial weights $\bf{W}$, outcome network $h_{\bf{V}}$ with initial weights $\bf{V}$, function family $\cF$ for IPM.
\STATE Compute $u=\frac{1}{n}\sum_{i=1}^n t_i$
 \STATE Compute $w_i = \frac{t_i}{2u} + \frac{1-t_i}{2(1-u)}$ for $i=1\ldots n$
  \WHILE{not converged}
   % \STATE Sample a mini-batch $\{(x_{i_j},t_{i_j},y_{i_j})\}_{j=1}^m$
       \STATE Sample mini-batch  $\{i_1,i_2,\ldots, i_m\} \subset  \{1,2,\ldots, n\}$
    \STATE Calculate the gradient of the IPM term:\\
    $g_1 = ${\small $\nabla_{\bf{W}} \; \text{IPM}_{\cF}(\{\Phi_{\bf{W}}(x_{i_j})\}_{t_{i_j} = 0}, \{\Phi_{\bf{W}}(x_{i_{k}})\}_{t_{i_j} = 1} )$}
    \STATE Calculate the gradients of the empirical loss: \\
    $g_2 = \nabla_{\bf{V}}\frac{1}{m} \sum_j w_{i_j} \cdot L\left(h_{\bf{V}}(\Phi_{\bf{W}}(x_{i_j}),t_{i_j}), y_{i_j}\right)$ \\
    $g_3 = \nabla_{\bf{W}}\frac{1}{m}\sum_j w_{i_j} \cdot L\left(h_{\bf{V}}(\Phi_{\bf{W}}(x_{i_j}),t_{i_j}), y_{i_j}\right)$
    \STATE Obtain step size scalar or matrix $\eta$ with standard neural net methods e.g. Adam~\citep{kingma2014adam}
    %\STATE Update $\mathbf{W} \leftarrow \mathbf{W} - \eta(\alpha g_1 +  g_3 )$
    %\STATE Update $\mathbf{V} \leftarrow \mathbf{V} - \eta(g_2 + 2 \lambda \mathbf{V})$
    \STATE  $\left[\mathbf{W},\mathbf{V}\right] \leftarrow \left[\mathbf{W} - \eta(\alpha g_1 +  g_3 ),  \mathbf{V} - \eta(g_2 + 2 \lambda \mathbf{V}) \right]$
    \STATE Check convergence criterion
  \ENDWHILE
\end{algorithmic}
\end{algorithm}
